\chapter{Visibility}
\label{ch:visibility}

\section{Visibility Problems}
\label{sec:visibility_problems}

The geometric notion of \textbf{visibility}\index{visibility} formalizes the
everyday question ``what can be seen from here?'': two points see each other
when the straight segment joining them is not blocked by intervening
obstacles. Visibility is at the heart of a surprising fraction of
computational geometry. Guarding an art gallery, rendering a scene, guiding
a robot, and finding the shortest route around a building all reduce, in
one way or another, to questions about which points see which other points.

\begin{definition}[Mutual Visibility]\label{def:vis_mutual}
  Let $\mathcal{O}$ be a finite set of pairwise-disjoint closed polygonal
  obstacles, or a simple polygon $P$ regarded as the only obstacle (so the
  free space is $P$ itself). Two points $p$ and $q$ that do not lie in the
  interior of an obstacle are \textbf{mutually visible} if the open segment
  between them does not intersect the interior of any obstacle; for a
  polygon, this is $\overline{pq} \subseteq P$.
\end{definition}

Visibility comes in several strengths, because the objects whose visibility
is of interest are usually segments, edges, or regions rather than bare
points.

\begin{definition}[Point, Edge, and Strong Visibility]\label{def:vis_edge}
  A point $p$ \textbf{sees} a segment $e$ if it sees at least one point of
  $e$. The point $p$ \textbf{strongly} (or \textbf{fully}) sees $e$ if it
  sees every point of $e$. Two segments $e_1$ and $e_2$ are mutually visible
  if some point of $e_1$ sees some point of $e_2$, and strongly visible if
  every point of $e_1$ sees every point of $e_2$.
\end{definition}

Guarding uses point visibility (every polygon point must be seen by a
guard), rendering needs the first surface hit along each viewing ray, and
shortest paths use edge visibility between obstacle vertices. The chapter
develops these in turn:

\begin{itemize}
  \item \textbf{Visibility inside polygons} (Section~\ref{sec:visibility_in_polygons}):
    the visibility relation within a simple polygon and its kernel;
  \item \textbf{Visibility polygons} (Sections~\ref{sec:visibility_polygons}
    --\ref{sec:visibility_from_point}): the region seen from a point and the
    algorithms that compute it;
  \item \textbf{Guarding} (Section~\ref{sec:visibility_guarding}): how many
    points suffice to see all of a polygon (the art gallery theorem);
  \item \textbf{Visibility graphs} (Section~\ref{sec:visibility_graphs}): the
    graph of mutually visible obstacle vertices, its complexity, and its
    computation;
  \item \textbf{Shortest and geodesic paths} (Sections~\ref{sec:visibility_shortest_paths}
    --\ref{sec:visibility_geodesic_paths}): visible-edge chains in the plane
    and inside polygons;
  \item \textbf{Applications} (Section~\ref{sec:visibility_applications}):
    hidden surface removal and visibility roadmaps.
\end{itemize}

The chapter connects to the rest of the book: the art gallery theorem is
developed fully in Chapter~\ref{ch:art_gallery}, visibility roadmaps are a
basic tool of Chapter~\ref{ch:motion_planning}, the rendering algorithms
reuse the machinery of Chapter~\ref{ch:applications_graphics}, the
envelopes of Chapter~\ref{ch:envelopes} describe a surface's silhouette
from a viewpoint, and the shortest-path results sit beside those of
Chapter~\ref{ch:distance_paths}.

\section{Visibility inside Polygons}
\label{sec:visibility_in_polygons}

Fix a simple polygon $P$ with $n$ vertices, its boundary oriented
counterclockwise. Inside $P$, visibility is the relation of
Definition~\ref{def:vis_mutual}: two points are visible when their
connecting segment lies in $P$. The relation is reflexive and symmetric but
\emph{not} transitive --- a point around a corner may see two points that do
not see each other --- which is why visibility is geometrically interesting
and algorithmically nontrivial.

A central structural object is the set of points from which \emph{everything}
is visible.

\begin{definition}[Kernel]\label{def:vis_kernel}
  The \textbf{kernel} $\operatorname{ker}(P)$ of a simple polygon $P$ is the
  set of points of $P$ that see every point of $P$. A polygon whose kernel is
  nonempty is called \textbf{star-shaped}.
\end{definition}

\begin{theorem}[Kernel Structure]\label{thm:vis_kernel}
  The kernel of a simple polygon $P$ is the intersection of the closed
  half-planes to the left of every directed edge of $P$ (edges oriented
  counterclockwise). It can be computed in $O(n)$ time, and $P$ is
  star-shaped if and only if its kernel is nonempty.
\end{theorem}

\begin{proof}
  A point $p$ sees an edge $e$ of $P$ if and only if $p$ lies in the
  half-plane determined by $e$ that is on the interior side of $P$; if $p$
  saw the interior side of every edge it would see the entire polygon, since
  $P$ is the intersection of those half-planes locally along the boundary.
  The kernel is therefore the intersection of the $n$ interior half-planes.
  Intersecting half-planes in the plane takes $O(n \log n)$ time in general;
  for a polygon the linear-time kernel algorithm of Lee and Preparata
  \cite{lee1979} exploits the cyclic order of the edges and runs in $O(n)$.
\end{proof}

The kernel is the ``point-visibility'' extremum: it is the largest set of
guards of size one, and star-shaped polygons are exactly the polygons that a
single guard can cover. The art gallery theorem of
Section~\ref{sec:visibility_guarding} asks how many points are needed when
one is not enough.

Between the kernel (one point) and all points lies the vertex-pair relation:
two vertices of $P$ are visible if the segment between them is a
\textbf{diagonal} of $P$, lying inside $P$ and meeting the boundary only at
its endpoints. Diagonals are the edges of every triangulation, the backbone
of visibility computation, because visibility questions inside a polygon
can be answered by walking across its triangles.

\begin{theorem}[Triangulation]\label{thm:vis_triangulation}
  Every simple polygon with $n$ vertices has a triangulation with exactly
  $n-2$ triangles and $n-3$ diagonals. Such a triangulation can be computed
  in $O(n \log n)$ time by the monotone-decomposition method
  (Chapter~\ref{ch:triangulations}) and in linear time by Chazelle's
  algorithm \cite{chazelle1991}.
\end{theorem}

Testing whether two vertices are mutually visible reduces to checking that
the segment between them is a diagonal: it must not cross any boundary edge
(the segment-intersection machinery of Chapter~\ref{ch:segment_intersection})
and must locally lie inside $P$ at both endpoints, a test using the
orientation predicates of Chapter~\ref{ch:geometric_predicates}. A simple
$O(n)$-per-pair check follows, so the visibility relation on the vertices is
computable in $O(n^2)$ after triangulation.

\index{visibility}\index{kernel}\index{star-shaped polygon}

\section{Visibility Polygons}
\label{sec:visibility_polygons}

\begin{definition}[Visibility Polygon]\label{def:vis_visibility_polygon}
  Let $p$ be a point of a simple polygon $P$ (or a point in the free space
  bounded by a set of obstacles). The \textbf{visibility polygon} of $p$,
  written $\operatorname{VP}(p)$, is the set of points $q$ that $p$ sees:
  \[
    \operatorname{VP}(p) = \{ q \in P : \overline{pq} \subseteq P \}.
  \]
\end{definition}

The visibility polygon is the exact region a viewer sees and a guard covers;
it is the query object underlying the guarding and navigation problems of
the next sections, and its structure is highly constrained.

\begin{theorem}[Structure of the Visibility Polygon]\label{thm:vis_vp_structure}
  Let $P$ be a simple polygon and $p \in P$ a point. Then
  $\operatorname{VP}(p)$ is a simple polygon, star-shaped with $p$ in its
  kernel. Its boundary consists of two kinds of edges:
  \begin{itemize}
    \item \textbf{boundary pieces}: maximal portions of edges of $P$ that
      are visible from $p$;
    \item \textbf{windows}: segments not lying on the boundary of $P$.
  \end{itemize}
  Every window is an open segment whose inner endpoint is a reflex vertex
  $v$ of $P$, which lies on the ray from $p$ through $v$, and which extends
  from $v$ to the first point of the boundary of $P$ met by that ray.
\end{theorem}

\begin{proof}
  The set $\operatorname{VP}(p)$ is star-shaped by definition: for every
  $q \in \operatorname{VP}(p)$ the segment $\overline{pq}$ is contained in
  $\operatorname{VP}(p)$, so $p$ lies in its kernel. The boundary of the
  region seen from $p$ consists of points at which the segment from $p$
  just grazes an obstacle. Along an edge of $P$ the visible part is a
  contiguous interval, since the boundary between two visible points of one
  edge is straight and unblocked. Where the visible boundary leaves
  $\partial P$, it does so at a reflex vertex $v$ (at a convex vertex the
  ray from $p$ continues to hug the boundary); the departing edge lies on
  the ray from $p$ through $v$ and extends until that ray hits $\partial P$
  again, beyond which $\overline{pq}$ is blocked. These are exactly the
  window edges of the statement.
\end{proof}

\begin{remark}\label{rem:vis_holes}
  If the obstacles are polygons with holes, or arbitrary disjoint segments,
  then $\operatorname{VP}(p)$ remains a star-shaped (hence simply connected)
  region: every visible point is joined to $p$ by a segment inside the
  region. The boundary now includes windows through reflex vertices of the
  obstacle polygons, and also segments joining the boundary of a hole to
  itself, but the star-shaped structure and the window characterization
  persist.
\end{remark}

\begin{example}[Visibility in a Staircase Polygon]\label{exa:vis_staircase}
  In a rectangular ``staircase'' of two steps, place the viewer $p$ at the
  bottom-left corner of the lower step. The visibility polygon contains the
  two horizontal boundary segments directly ahead, the vertical walls
  opposite, and the portion of the upper step's floor not occluded by the
  lower step's wall. The occlusion ends at a window: a segment through the
  reflex corner atop that wall, running along the ray from $p$ through that
  corner until it hits the far wall, beyond which the upper floor is visible
  again --- the prototypical ``see around a corner.''
\end{example}

\index{visibility polygon}\index{window}

\section{Computing Visibility from a Point}
\label{sec:visibility_from_point}

\subsection{1. Overview}

Two algorithms dominate the computation of $\operatorname{VP}(p)$:
\begin{itemize}
  \item \textbf{Lee's algorithm} \cite{lee1983} computes
    $\operatorname{VP}(p)$ for a simple polygon in $O(n)$ time by processing
    the vertices in the rotational order induced by the polygon and
    maintaining a stack of currently visible vertices; it is the historical
    first optimal algorithm.
  \item \textbf{The rotational sweep} sweeps a ray around $p$ and records
    the closest obstacle along each direction. It runs in $O(n \log n)$
    time and, unlike Lee's algorithm, generalizes to polygons with holes and
    to arbitrary sets of disjoint segments.
\end{itemize}
Guibas, Hershberger, Leven, Sharir, and Tarjan \cite{guibas1987vis} showed
that after triangulation the visibility polygons of \emph{all} vertices of a
simple polygon can be computed in $O(n)$ total time, and Laszlo
\cite{laszlo1996} gives an implementer-oriented account. This section
presents the rotational sweep and how Lee's algorithm sharpens it to linear
time for simple polygons.

\subsection{2. Definitions}

\begin{definition}[Radial Order]\label{def:vis_radial_order}
  The \textbf{radial order} of a set of points around $p$ is their cyclic
  order by polar angle around $p$. Two points at the same angle are ordered
  by distance from $p$.
\end{definition}

\begin{definition}[Front]\label{def:vis_front}
  During a sweep of a ray rotating about $p$, the \textbf{front} is the
  obstacle segment that intersects the ray at the smallest distance from $p$,
  i.e., the segment currently occluding the view along the ray. The front
  changes only at \textbf{events}, the angles at which the ray passes an
  endpoint of an obstacle segment.
\end{definition}

The window edges of Theorem~\ref{thm:vis_vp_structure} are the records of
front changes: when the front jumps discontinuously from one segment to
another at a reflex vertex, the jump is a window.

\subsection{3. Theory}

The rotational sweep is correct because visibility is radial: $q$ is
visible from $p$ if and only if no obstacle lies strictly between $p$ and
$q$ on the ray through $q$, so in every angular direction the boundary of
$\operatorname{VP}(p)$ is the closest obstacle point. Between consecutive
event angles the set of segments crossed by the ray is fixed, so the visible
boundary is one straight piece of the front, clipped to that interval. At an
event the front may change; if the hit point jumps from $a$ to $b$, the
segment $\overline{ab}$ is a window (through a reflex vertex, by
Theorem~\ref{thm:vis_vp_structure}). Sweeping all $O(n)$ events and emitting
the visible pieces and windows therefore yields exactly
$\operatorname{VP}(p)$.

Lee's linear-time algorithm \cite{lee1983} removes the sorting step using
two structural facts about simple polygons: (i) the radial order of the
vertices around an interior point is obtainable in linear time without
sorting, because the polygon's edges restrict how that order can change; and
(ii) each polygon edge contributes only a bounded number of changes to the
stack of visible vertices. Processing vertices in that order with a stack
reproduces the sweep's output in $O(n)$ total time: the correctness transfers
from the sweep, and only the bound is sharpened.

\subsection{4. Pseudocode}

\begin{algorithm}[htbp]
\caption{Rotational Sweep for the Visibility Polygon}
\label{alg:vis_rotational_sweep}
\begin{algorithmic}[1]
\Require A point $p$ and the edges of a simple polygon (or of disjoint obstacle polygons), in general position.
\Ensure The visibility polygon $\operatorname{VP}(p)$ as a cyclic sequence of boundary points.
\Function{VP\_Sweep}{$p$, edges of the scene}
  \State $E \gets$ all edge endpoints sorted by angle around $p$ (ties by distance)
  \State $\ell \gets$ the front segment at the initial sweep angle
  \State $B \gets \emptyset$ \Comment{visible boundary, emitted in angular order}
  \For{event vertex $v \in E$, in angular order}
    \State $\theta \gets$ angle of $v$
    \State $a \gets$ hit point of the front $\ell$ with the ray at angle $\theta$
    \State update the front: insert the segment(s) incident to $v$, delete the segment whose angular interval ended
    \State $b \gets$ hit point of the updated front with the ray at angle $\theta$
    \If{$a \neq b$}
      \State emit the window $\overline{ab}$ (passing through a reflex vertex)
    \EndIf
    \State emit the visible piece of the front between angle $\theta$ and the next event
  \EndFor
  \State emit the final front piece and close the polygon
  \State \Return $B$
\EndFunction
\end{algorithmic}
\end{algorithm}

\subsection{5. Parameter Selection}

\begin{itemize}
\item \textbf{General position}: the sweep assumes no two events coincide
  and no event lies exactly at the sweep start angle. Coincident angles are
  ordered by distance; the case of the ray passing exactly through a vertex
  is resolved by a fixed convention (treat the vertex as reached by the ray,
  then reprocessed after updating the front).
\item \textbf{Tie-breaking}: when several vertices share an angle, they must
  be processed in distance order from $p$; the closest reflex vertex is the
  one that can create a window.
\item \textbf{Robust predicates}: all front updates and hit-point
  computations are orientation and intersection tests; they must use the
  exact predicates of Chapter~\ref{ch:geometric_predicates} to avoid
  inconsistent front states.
\end{itemize}

\subsection{6. Complexity}

\begin{itemize}
\item \textbf{Time}: $O(n \log n)$ for the rotational sweep, dominated by
  the radial sort of $n$ endpoints; $O(n)$ for Lee's algorithm on a simple
  polygon \cite{lee1983}.
\item \textbf{Space}: $O(n)$ for both algorithms (the sorted event list and
  the output polygon).
\item \textbf{Lower bound}: $\Omega(n)$ trivially, since the output may have
  $\Theta(n)$ vertices, so Lee's algorithm is optimal for simple polygons.
\end{itemize}

\subsection{7. Usages}

\begin{itemize}
\item \textbf{Guarding}: $\operatorname{VP}(p)$ is the region covered by a
  guard at $p$ (Section~\ref{sec:visibility_guarding}).
\item \textbf{Visibility graph}: computing the visibility graph of a
  polygon's vertices reduces to computing $\operatorname{VP}(v)$ for every
  vertex $v$ (Section~\ref{sec:visibility_graphs}); the vertices visible
  from $v$ are exactly the polygon vertices on the boundary of
  $\operatorname{VP}(v)$.
\item \textbf{Sensor models}: in robotics and games, the visibility polygon
  is the exact field of view of a sensor or camera, used for coverage,
  surveillance, and line-of-sight queries.
\end{itemize}

\subsection{8. Testing Methods and Edge Cases}

\begin{itemize}
\item \textbf{Brute-force reference}: sample points densely in a bounding
  box and compare against a brute-force visible test; every sampled visible
  point must lie in the computed $\operatorname{VP}(p)$ and conversely.
\item \textbf{Star-shaped check}: verify that $p$ lies in the kernel of the
  computed polygon (Section~\ref{sec:visibility_in_polygons}) and that no
  two edges of the output cross.
\item \textbf{Window audit}: every non-boundary edge of the output must be a
  segment through a reflex vertex of $P$ on the ray from $p$
  (Theorem~\ref{thm:vis_vp_structure}); the number of windows agrees with
  the number of such reflex vertices up to $O(1)$.
\item \textbf{Degenerate inputs}: points on the boundary of $P$ (where
  $\operatorname{VP}(p)$ degenerates to a wedge), collinear vertices, and
  nested holes must be exercised.
\end{itemize}

\subsection{9. References}

Lee's linear-time algorithm is the classic reference \cite{lee1983};
the linear-time treatment of visibility and shortest paths from a
triangulation, together with the radial-order lemma for simple polygons, is
in \cite{guibas1987vis}. Textbook presentations appear in \cite{orourke1987,
berg2008,laszlo1996}, and the robust implementation issues follow
Chapter~\ref{ch:geometric_predicates}.

\index{rotational sweep}\index{Lee's algorithm}

\section{Guarding Polygons}
\label{sec:visibility_guarding}

The \textbf{art gallery problem}\index{art gallery problem} asks: how many
points (guards) inside a polygon are needed so that every point of the
polygon is seen by at least one of them? A guard at $p$ covers exactly
$\operatorname{VP}(p)$ (Section~\ref{sec:visibility_polygons}), so the
question is a covering problem for the family of visibility polygons.

\begin{theorem}[Art Gallery Theorem]\label{thm:vis_art_gallery}
  Every simple polygon with $n$ vertices can be guarded by at most
  $\lfloor n/3 \rfloor$ guards, and this bound is tight: for every $n$
  there are polygons that cannot be guarded by fewer than $\lfloor n/3
  \rfloor$ guards.
\end{theorem}

\begin{proof}
  The upper bound is Fisk's elegant proof \cite{fisk1978}. Triangulate the
  polygon (Theorem~\ref{thm:vis_triangulation}). The triangulation graph is
  3-colorable: its dual graph is a tree, so one can 3-color the vertices by
  walking the dual tree, each new triangle forced to receive the third color
  on its new vertex. Of the three color classes, one contains at most
  $\lfloor n/3 \rfloor$ vertices. Place guards at the vertices of that
  class. Every triangle has exactly one vertex of each color, so every
  triangle contains a guard at one of its vertices; since a triangle is
  convex, that guard sees the entire triangle, and every point of the
  polygon lies in some triangle. The polygon is therefore guarded by at most
  $\lfloor n/3 \rfloor$ guards.
\end{proof}

The lower bound is a \textbf{comb polygon}\index{comb polygon}: a long thin
rectangle with $\lfloor n/3 \rfloor$ triangular ``teeth'' cut into one side
so that the teeth are separated by narrow necks. A single guard can see
points in at most one tooth plus the adjacent corridor, so at least one
guard per tooth is required; each tooth costs three vertices, giving the
$\lfloor n/3 \rfloor$ lower bound. The theorem is Chvatal's
\cite{chvatal1975}; the full history and the algorithmic extensions are in
Chapter~\ref{ch:art_gallery}.

\begin{remark}\label{rem:vis_guarding_variants}
  Variants change the picture. \textbf{Vertex guards} (restricted to
  vertices) satisfy the same $\lfloor n/3 \rfloor$ bound, since Fisk's proof
  already produces vertex guards. For polygons with $h$ holes the bound is
  $\lfloor (n + 2h)/3 \rfloor$, and for orthogonal polygons $\lfloor n/4
  \rfloor$. Computing the \emph{minimum} number of guards is NP-hard even
  for simple polygons, so the worst-case bound is the practical substitute
  for exact optimization.
\end{remark}

The guarding problem is the point-coverage version of set cover, with the
visibility polygons of Section~\ref{sec:visibility_polygons} as the sets to
be covered; Chapter~\ref{ch:art_gallery} develops the theorem, its variants,
and the hardness results in depth.

\index{guard}\index{art gallery theorem}

\section{Visibility Graphs}
\label{sec:visibility_graphs}

\subsection{1. Overview}

The \textbf{visibility graph}\index{visibility graph} of a set of obstacles
records which obstacle vertices see each other. Its paths are exactly the
candidate shortest paths among the obstacles: a path that changes direction
only at obstacle corners moves in a straight line between consecutive
vertices precisely when those vertices are mutually visible. The visibility
graph therefore turns the continuous shortest-path problem into a discrete
graph problem (Section~\ref{sec:visibility_shortest_paths}). This section
studies the base case of pairwise-disjoint segments, as analyzed in
\cite{berg2008}.

\subsection{2. Definitions}

\begin{definition}[Visibility Graph]\label{def:vis_vg}
  Let $\mathcal{S}$ be a set of pairwise-disjoint closed line segments
  (the obstacles). The \textbf{visibility graph} of $\mathcal{S}$ has a
  vertex for every segment endpoint and an edge between two vertices $u, v$
  if $u$ and $v$ are mutually visible: the open segment $\overline{uv}$
  crosses no segment of $\mathcal{S}$. For polygonal obstacles one takes the
  obstacle vertices as vertices and requires the open segment to avoid all
  obstacle interiors.
\end{definition}

\begin{definition}[Bitangent]\label{def:vis_bitangent}
  A \textbf{bitangent} (or visibility edge) of a set of disjoint convex
  obstacles is a segment that is mutually visible to the two obstacles it
  touches at its endpoints. In the segment-obstacle setting the bitangents
  are exactly the visibility-graph edges, since every visibility edge is
  tangent to the segments at both endpoints.
\end{definition}

\subsection{3. Theory}

\begin{theorem}[Complexity of the Visibility Graph]\label{thm:vis_complexity}
  The visibility graph of $n$ pairwise-disjoint segments has $O(n^2)$ edges,
  and the bound is tight: configurations with $O(n^2)$ mutual-visibility
  pairs exist. The graph can be computed in $O(n^2 \log n)$ time by a
  per-vertex rotational sweep, in $O(n^2)$ time with more care, and in
  $O(n \log n + k)$ time with an output-sensitive algorithm, where $k$ is
  the number of visibility edges \cite{ghoshmount1991,overmarswelzl1992}.
\end{theorem}

The rotational sweep analyzes visibility from one vertex $v$ at a time.
Sort the other endpoints around $v$ and rotate a ray, maintaining the front
(Definition~\ref{def:vis_front}), the closest segment crossing the ray. When
the ray points at a vertex $u$, the open segment $\overline{vu}$ is crossed
by an obstacle if and only if the front intersects it: a closer blocker
would be the front, and a farther blocker would be beyond it. Hence $u$ is
visible from $v$ exactly when the front does not cross $\overline{vu}$. The
sweep costs $O(n \log n)$ per vertex, giving $O(n^2 \log n)$; the
output-sensitive methods of \cite{ghoshmount1991,overmarswelzl1992} spend
time proportional to the number of edges actually reported.

For a simple polygon, the visibility graph is obtained by computing
$\operatorname{VP}(v)$ for every vertex $v$
(Section~\ref{sec:visibility_from_point}): the vertices visible from $v$
are exactly the polygon vertices on the boundary of $\operatorname{VP}(v)$.
Since each visibility polygon has $O(n)$ boundary vertices, all $n$
computations take $O(n^2)$ time, matching the worst-case number of edges.
The bitangent structure connects to the envelopes of Chapter~\ref{ch:envelopes}:
the visibility edges from a vertex are the vertices of the lower envelope of
the obstacle segments as seen from that vertex.

\subsection{4. Pseudocode}

\begin{algorithm}[htbp]
\caption{Rotational Sweep for the Visibility Graph}
\label{alg:vis_vg_sweep}
\begin{algorithmic}[1]
\Require Pairwise-disjoint obstacle segments $\mathcal{S}$ in general position.
\Ensure The visibility graph $G = (V, E)$.
\Function{VisibilityGraph}{$\mathcal{S}$}
  \State $V \gets$ all segment endpoints; $E \gets \emptyset$
  \For{vertex $v \in V$}
    \State $W \gets$ the other endpoints of $V$, sorted by angle around $v$
    \State $\ell \gets$ the front segment at the initial sweep angle from $v$
    \For{endpoint $u \in W$, in radial order}
      \If{the front $\ell$ does not cross the open segment $\overline{vu}$}
        \State $E \gets E \cup \{(v,u)\}$ \Comment{$u$ is visible from $v$}
      \EndIf
      \State update $\ell$ as the ray passes the endpoints of the segments incident to $u$
    \EndFor
  \EndFor
  \State \Return $G$
\EndFunction
\end{algorithmic}
\end{algorithm}

\subsection{5. Parameter Selection}

\begin{itemize}
\item \textbf{General position}: assume no two endpoints share an angle from
  $v$ and no endpoint lies on another segment; ties and collinearities are
  resolved by the fixed conventions of Chapter~\ref{ch:geometric_predicates}.
\item \textbf{Distinct obstacles}: input segments must be pairwise disjoint;
  touching segments are split at their contact points first.
\item \textbf{Output size}: the graph may be dense; applications should
  build only the edges they need (e.g., the reduced visibility graph keeping
  only bitangent edges between reflex vertices plus tangent edges to convex
  vertices) when $k$ is large.
\end{itemize}

\subsection{6. Complexity}

\begin{itemize}
\item \textbf{Time}: $O(n^2 \log n)$ for Algorithm~\ref{alg:vis_vg_sweep}
  ($O(n \log n)$ per vertex); $O(n \log n + k)$ output-sensitive
  \cite{ghoshmount1991,overmarswelzl1992}.
\item \textbf{Space}: $O(n + k)$ for the graph and auxiliary structures.
\item \textbf{Lower bound}: $\Omega(n^2)$ time in the worst case, since the
  output can have $\Theta(n^2)$ edges.
\end{itemize}

\subsection{7. Usages}

\begin{itemize}
\item \textbf{Shortest paths}: the shortest obstacle-avoiding path is a
  visibility-graph path (Section~\ref{sec:visibility_shortest_paths}).
\item \textbf{Motion planning}: visibility graphs are the classical
  \textbf{roadmaps}\index{roadmap} for a point robot moving among polygonal
  obstacles; Chapter~\ref{ch:motion_planning} uses them as the simplest
  complete roadmap.
\item \textbf{Navigation meshes}: reduced visibility graphs (bitangents
  between reflex vertices) form the skeleton of navigation meshes for agents
  in games and robotics.
\end{itemize}

\subsection{8. Testing Methods and Edge Cases}

\begin{itemize}
\item \textbf{Brute-force verification}: for small $n$, compute all pairs
  $(u,v)$ by a direct segment-intersection test against every obstacle and
  require the same edge set.
\item \textbf{Edge validity}: every reported edge must be visible, and no
  omitted pair may be visible, on random and on adversarial inputs.
\item \textbf{Tangency}: endpoints lying on the same ray from $v$, and
  segments whose endpoints are collinear, must produce a consistent edge set
  under the fixed tie-breaking rule.
\item \textbf{Dense instances}: a ring of short segments with a central
  vertex should produce a near-complete visibility graph; verify the count
  against $\binom{2n}{2}$.
\end{itemize}

\subsection{9. References}

The visibility graph and its rotational-sweep computation are presented in
\cite{berg2008}; the optimal output-sensitive algorithms are due to Ghosh
and Mount \cite{ghoshmount1991} and to Overmars and Welzl
\cite{overmarswelzl1992}; the use of visibility graphs in planning dates to
the configuration-space approach of Lozano-Perez \cite{lozanoperez1983} and
is treated in \cite{orourke1987,sharir1995}.

\section{Shortest Paths among Obstacles}
\label{sec:visibility_shortest_paths}

The shortest-path problem among obstacles is: given two points $s, t$ and
pairwise-disjoint polygonal obstacles, find the shortest path from $s$ to
$t$ avoiding the obstacle interiors --- the "rubber band" path pulled tight
between $s$ and $t$. Its combinatorial structure is the central fact that
makes visibility graphs useful.

\begin{theorem}[Shortest Path Structure]\label{thm:vis_shortest_path}
  Let $s, t$ be points and $\mathcal{O}$ a set of pairwise-disjoint polygonal
  obstacles. The shortest obstacle-avoiding path from $s$ to $t$ is a simple
  polygonal chain; every interior vertex of the chain is a vertex of an
  obstacle, and every two consecutive vertices of the chain are mutually
  visible. Consequently the shortest path is a path in the visibility graph
  of the obstacle vertices augmented with $s$ and $t$.
\end{theorem}

\begin{proof}
  A shortest path avoids obstacle interiors. Any subpath that runs through
  free space with no obstacle corner between its endpoints can be replaced
  by the straight segment between its endpoints, which shortens it; hence
  every ``kink'' occurs at an obstacle vertex, and the path is a polygonal
  chain through obstacle vertices. If two consecutive vertices $u, v$ of the
  chain were not mutually visible, $\overline{uv}$ would cross an obstacle
  and the path could be locally detoured to a shorter one, contradicting
  optimality. Hence consecutive vertices are visible, and the chain is a
  path in the visibility graph of the vertices plus $s$ and $t$.
\end{proof}

The theorem reduces the continuous problem to a shortest-path problem in a
graph with $O(n)$ vertices and $O(n^2)$ edges:

\begin{algorithm}[htbp]
\caption{Shortest Path among Polygonal Obstacles}
\label{alg:vis_shortest_path}
\begin{algorithmic}[1]
\Require Points $s, t$ and pairwise-disjoint polygonal obstacles $\mathcal{O}$.
\Ensure The shortest obstacle-avoiding path from $s$ to $t$.
\Function{ShortestPath}{$s, t, \mathcal{O}$}
  \State $G \gets$ visibility graph of $\mathcal{O}$ (Algorithm~\ref{alg:vis_vg_sweep})
  \State add $s$ and $t$ to $G$ with their visible edges
  \State run Dijkstra's algorithm on $G$ with edge weights equal to Euclidean lengths
  \State \Return the resulting $s$-$t$ path
\EndFunction
\end{algorithmic}
\end{algorithm}

With $O(n^2)$ edges the Dijkstra run costs $O(n^2 \log n)$; with the
output-sensitive visibility graph of Theorem~\ref{thm:vis_complexity} the
total is $O(n \log n + k)$. The general Euclidean shortest-path problem and
its rectilinear and weighted variants are discussed in
Chapter~\ref{ch:distance_paths}; the visibility graph remains the reference
approach whenever the shortest path bends only at obstacle corners
\cite{lozanoperez1983,berg2008}.

\begin{remark}\label{rem:vis_spp_directions}
  Three refinements are standard. (i) The \textbf{reduced} visibility graph
  (bitangent edges between reflex vertices plus tangents to convex vertices)
  still contains the shortest path and is much sparser in practice. (ii) For
  a fixed start and many targets, the graph is processed once and each query
  is a Dijkstra-like exploration. (iii) For a robot of nonzero size, the
  obstacles are first inflated to a configuration space
  (Section~\ref{sec:visibility_applications}).
\end{remark}

\index{shortest path}

\section{Geodesic Paths}
\label{sec:visibility_geodesic_paths}

When the path is required to stay inside a simple polygon, the shortest
path is a \textbf{geodesic path}\index{geodesic path}: the shortest curve
in $P$ joining two points of $P$. The geodesic is the planar analogue of
the great-circle distance on a surface, and it defines the \textbf{geodesic
distance} $d_P(s,t)$, the length of the geodesic.

\begin{definition}[Geodesic]\label{def:vis_geodesic}
  For points $s, t$ in a simple polygon $P$, the \textbf{geodesic} $\pi(s,t)$
  is the shortest path from $s$ to $t$ whose points all lie in $P$, and
  $d_P(s,t)$ is its length.
\end{definition}

\begin{theorem}[Structure and Computation of Geodesics]\label{thm:vis_geodesic}
  The geodesic $\pi(s,t)$ is unique and simple, and all its interior
  vertices are reflex vertices of $P$. Given a triangulation of $P$, the
  geodesic can be computed in $O(n)$ time.
\end{theorem}

\begin{proof}
  Uniqueness follows because the free space is simply connected: any two
  distinct shortest paths of equal length between the same endpoints enclose
  a region of $P$ of positive area, and a slight shortening within that
  region contradicts optimality. As in
  Theorem~\ref{thm:vis_shortest_path}, every bend of the geodesic must lie
  at a reflex vertex of $P$: at an interior bend the path could be shortcut,
  and at a convex boundary vertex the path could be pulled into the interior
  of $P$. The geodesic crosses exactly the triangles on the unique path
  between the triangles of $s$ and $t$ in the dual tree of the triangulation,
  and the linear-time algorithms of \cite{guibas1987vis} walk this chain
  while maintaining a \textbf{funnel}: the shortest path to the current
  boundary of the visited region, which has the shape of a funnel with the
  source $s$ at the apex.
\end{proof}

The funnel algorithm processes the triangle chain one at a time. Each step
adds a triangle contributing one new vertex; the funnel's two boundary
chains (the shortest paths from $s$ to the two ends of the path front) are
updated by a visibility test of which vertex lies inside the current
triangle. When the two chains cross, the apex advances to the crossing point
and the funnel resets --- the geodesic's turning point moves forward. Each
triangle is processed once, so the total cost is linear after locating $s$
and $t$ in the triangulation (Chapter~\ref{ch:point_location}).

\begin{algorithm}[htbp]
\caption{Shortest Path in a Triangulated Polygon (Funnel Algorithm)}
\label{alg:vis_funnel}
\begin{algorithmic}[1]
\Require A simple polygon $P$ with a triangulation $T$, and points $s, t$ in $P$.
\Ensure The geodesic path from $s$ to $t$.
\Function{FunnelShortestPath}{$P, T, s, t$}
  \State locate the triangles $\Delta_s \ni s$ and $\Delta_t \ni t$ in $T$
  \State $\Delta_s, \Delta_1, \dots, \Delta_k, \Delta_t \gets$ the triangle chain along the dual tree of $T$
  \State $F \gets$ funnel with apex $s$ and the two boundary chains along the current path front
  \For{triangle $\Delta$ in the chain}
    \State add the new vertex of $\Delta$; update the two funnel chains
    \If{the two chains cross at a point $c$}
      \State advance the apex to $c$ and reset the funnel
    \EndIf
  \EndFor
  \State \Return the shortest path read off the final funnel
\EndFunction
\end{algorithmic}
\end{algorithm}

\begin{example}[Two Geodesics Bound a Funnel]\label{exa:vis_funnel}
  Fix a source $s$ and a reflex vertex $v$ of $P$. The two geodesics
  $\pi(s,v)$ together with the boundary of $P$ between them enclose a
  funnel-shaped region; every geodesic from $s$ to a point $t$ in that
  region follows one of the two paths until the funnel narrows, then
  crosses to the other side. This structure is the organizing principle of
  \cite{guibas1987vis} and of data structures answering geodesic-distance
  queries in logarithmic time after $O(n \log n)$ preprocessing.
\end{example}

The geodesic notion extends to polyhedral surfaces, where a shortest path
is a sequence of straight segments across faces meeting only at surface-reflex
vertices. The \textbf{discrete geodesic problem} \cite{mitchell1987} computes
such paths by continuous Dijkstra simulation, together with the polygon case
in Chapter~\ref{ch:distance_paths}; visibility enters because the path's
straight segments must lie in a mutually visible region.

\index{funnel algorithm}\index{geodesic}

\section{Applications}
\label{sec:visibility_applications}

Visibility is a lens through which several mature application areas are best
understood. Three stand out: rendering, motion planning, and robotics.

\subsection*{Hidden Surface Removal and Rendering}

The \textbf{hidden surface removal} problem is to draw a 3D scene so that
only the surfaces visible from the viewpoint appear. It is, per pixel, a
point-visibility problem: the pixel's color comes from the \emph{first}
surface hit along the viewing ray through the pixel. The two classical
solutions are:

\begin{description}
  \item[The painter's algorithm] \cite{newell1972} sorts the polygons by
    depth and draws them back-to-front, letting near polygons overpaint far
    ones. It fails when no consistent depth order exists (cyclically
    overlapping polygons), requiring the polygons to be split.
  \item[The z-buffer] \cite{catmull1974} performs the depth comparison per
    pixel: a fragment is written only if its depth beats the stored depth,
    so the buffer holds exactly the per-pixel first hit at $O(1)$ per
    fragment --- the hardware-standard algorithm of modern graphics
    pipelines.
  \item[Binary space partitions] \cite{fuchs1980,naylor1990} build a
    recursive partition of space by planes; any viewpoint then has a
    correct back-to-front order from a single tree walk, giving the
    painter's algorithm its missing order.
\end{description}

Geometrically, the visible surface points are the \textbf{lower envelope}
of the surfaces as seen from the viewpoint: a surface point is visible
exactly when it lies on that envelope. This ties rendering to the envelope
machinery of Chapter~\ref{ch:envelopes} and, on the engineering side, to
the spatial structures of \cite{ericson2004} used for occlusion culling;
Chapter~\ref{ch:applications_graphics} uses these algorithms as its
rendering core.

\subsection*{Motion Planning and Visibility Roadmaps}

The configuration-space approach of Lozano-Perez \cite{lozanoperez1983}
reduces motion planning for a robot with geometry to point planning for its
\emph{configuration} among inflated obstacles (the \textbf{C-obstacles}).
For a point robot among polygonal obstacles the visibility graph is a
complete roadmap: every shortest path is a visibility-graph path
(Theorem~\ref{thm:vis_shortest_path}), so planning reduces to the
shortest-path search of Algorithm~\ref{alg:vis_shortest_path}. Visibility
roadmaps remain the reference ``exact'' planner for low-dimensional
configuration spaces, and visibility is also the connection criterion in
sampling-based planners. Chapter~\ref{ch:motion_planning} develops the
roadmap approach and its sampling-based successors.

\subsection*{Robotics, Coverage, and Simulation}

Guarding (Section~\ref{sec:visibility_guarding}) is the coverage question
behind sensor placement: placing cameras, guards, or beacons so that every
point of a domain is seen by at least one sensor, with the art gallery
bound as worst-case guarantee (optimal placement is NP-hard,
Remark~\ref{rem:vis_guarding_variants}). In navigation, the visibility
polygon is the local sensor model of simultaneous localization and mapping,
and the geodesic paths of Section~\ref{sec:visibility_geodesic_paths} are
the routes agents follow indoors. In games and simulation, reduced
visibility graphs form navigation meshes, and the point-visibility and
ray-shooting primitives of this chapter underlie the collision and occlusion
queries discussed in \cite{ericson2004}.

\index{hidden surface removal}\index{painter's algorithm}\index{z-buffer}
\index{binary space partition}\index{visibility roadmap}

\section{Exercises}
\label{sec:visibility_exercises}

\begin{enumerate}
  \item Show that mutual visibility is not transitive: construct three
    points $p, q, r$ in a simple polygon with $p$ and $q$ visible, $q$ and
    $r$ visible, but $p$ and $r$ not visible. Prove that in a convex polygon
    the relation \emph{is} transitive.
  \item Prove that the kernel of a simple polygon is convex and is exactly
    the intersection of the half-planes of Definition~\ref{def:vis_kernel},
    and give a polygon whose kernel is a single point.
  \item Prove Theorem~\ref{thm:vis_vp_structure} in detail: show that every
    window of $\operatorname{VP}(p)$ passes through a reflex vertex of $P$
    and lies on the ray from $p$ through that vertex.
  \item Carry out Fisk's proof of Theorem~\ref{thm:vis_art_gallery}: prove
    that the triangulation graph of a simple polygon is 3-colorable by
    induction on the dual tree, and derive the $\lfloor n/3 \rfloor$ bound.
  \item Design a comb polygon on $n$ vertices and prove that it cannot be
    guarded with fewer than $\lfloor n/3 \rfloor$ guards.
  \item Prove Theorem~\ref{thm:vis_shortest_path}: show that every kink of
    the shortest obstacle-avoiding path occurs at an obstacle vertex and
    that consecutive vertices of the path are mutually visible.
  \item Compute by hand the visibility polygon of a point in a staircase
    polygon, and verify that the number of windows equals the number of
    reflex vertices visible on windows, as guaranteed by
    Theorem~\ref{thm:vis_vp_structure}.
  \item Implement the rotational sweep of
    Algorithm~\ref{alg:vis_rotational_sweep} for a simple polygon and
    validate it against a brute-force visibility test on random polygons,
    checking in particular the window audit described in
    Section~\ref{sec:visibility_from_point}.
\end{enumerate}
